\chapter{Probability Semantics at \texorpdfstring{$c_*$}{c*}}
\label{ch:probability}

On which measurable objects, and in which order, does every probabilistic statement at $c_*$ rest? Laws
and kernels come first, densities follow as Radon--Nikodym derivatives against references fixed for a
whole family, and logarithms of density ratios enter only under
\Cref{def:prob-log-density-ratio-convention}; this order makes the evidence bound of \Cref{ch:elbo} an
identity in the extended reals. The main results are common domination of a mixed family exactly when its
atoms are countable (\Cref{thm:prob-mixed-family-common-domination}), the relative-entropy chain rule in
$[0,+\infty]$ (\Cref{lem:prob-kl-chain-rule}) that every later decomposition cites, and three levels of
description, at which a law-point energy is not a sample energy and point-mass recognition of law points
has infinite cost (\Cref{prop:prob-law-point-energy,prop:prob-law-point-point-mass}).
\Cref{ch:generative} builds the joint on these spaces, \Cref{ch:elbo} applies the chain rule, and
\Cref{ch:selective-communication,ch:finite-depth-formation} use the levels.

\paragraph{Scope.}
Injectivity of a model evaluator is not assumed, and no regularity of the quotient of presentations with
the same evaluated kernel is supplied; presentations are retained. The coupling selector is applied once,
and its equivariance under frame re-expression (rerunning it on transformed marginals returns the
transformed coupling) is neither assumed nor claimed. Open are a normalized law-point prior
(\Cref{open:prob-law-point-prior}) and likelihoods for record families without a common dominating
measure (\Cref{open:prob-nondominated-likelihood}). \status{NOT-CLAIMED}

\section{Latent and record spaces at \texorpdfstring{$c_*$}{c*}}
\label{sec:prob-design}

\definitionheading{Sample spaces and product sigma-algebras}{def:prob-product-sample-spaces}
For each agent $i\in V$ let $(\mathsf K_i,\mathscr K_i)$ and $(\mathsf M_i,\mathscr M_i)$ be its belief
and model sample fibers at $c_*$, identified with $\mathsf K$ and $\mathsf M$ by the frame $u_i$, and put
$\mathsf Y_i=\mathsf K_i\times\mathsf M_i$. Let $\mathfrak A$ be a finite set of records; each
$a\in\mathfrak A$ has a scope $\partial a\subseteq V$ and a record space $(\mathsf O_a,\mathscr O_a)$. The
latent and observation spaces are
\begin{align}
\mathsf Y&=\prod_{i\in V}\mathsf Y_i,
&\mathscr Y&=\bigotimes_{i\in V}(\mathscr K_i\otimes\mathscr M_i),
\label{eq:prob-latent-space}\\
\mathsf O&=\prod_{a\in\mathfrak A}\mathsf O_a,
&\mathscr O&=\bigotimes_{a\in\mathfrak A}\mathscr O_a,
\label{eq:prob-observation-space}
\end{align}
with points $y=((k_i,m_i))_{i\in V}$ and $o=(o_a)_{a\in\mathfrak A}$, and
$\mathsf Y_{\partial a}=\prod_{i\in\partial a}\mathsf Y_i$. An observation of a single agent is a record
with singleton scope.

A population sample is a point of $\mathsf Y$, and evaluating a sample field at $c_*$ returns
\begin{equation}
\operatorname{Ev}(y)=\bigl((k_i(c_*),m_i(c_*))\bigr)_{i\in V}\in\mathsf Y,
\label{eq:prob-evaluation-map}
\end{equation}
whereas the section values $q_i^b,q_i^m$ at $c_*$ are laws on $\mathsf K_i$ and $\mathsf M_i$, which become
latent coordinates only at the law-point level of \Cref{sec:prob-levels}. Missing or aggregated
observations are represented by an observation kernel into a coarser record space, which changes
$\mathsf O$ and leaves $\mathsf Y$ unchanged.

\section{Structural data and kernel signatures}
\label{sec:prob-structural}

\definitionheading{Structural configuration and kernel signatures}{def:prob-structural-kernel-signatures}
Fix a measurable configuration space $(\mathsf X,\mathscr X)$. Every $X\in\mathsf X$ uses the same agent
set $V$, the same records and scopes, and the same fiber identifications, so $(\mathsf Y,\mathscr Y)$ and
$(\mathsf O,\mathscr O)$ do not vary with $X$. A configuration carries the interaction multigraph
$\mathfrak I$ of \Cref{def:geo-graph-links}, its links $\Theta_e^x$, the frames $u_i$, and the parameters
of the mechanisms. With $(\mathsf I_i,\mathscr I_i)$ an access space for each $i\in V$, the access maps,
local recognition kernels, record kernels, and population laws have the signatures
\begin{align}
\operatorname{Acc}_i&:\mathsf O\times\mathsf X\longrightarrow\mathsf I_i,
\quad\mathbb Q_i:\mathsf I_i\rightsquigarrow\mathsf Y_i,
&\mathbb Q_{i,o,X}&:=\mathbb Q_i(\operatorname{Acc}_i(o,X),\cdot),
\label{eq:prob-local-recognition-kernel}\\
K_a&:(\mathsf X\times\mathsf Y_{\partial a})\rightsquigarrow\mathsf O_a,
&a&\in\mathfrak A,
\label{eq:prob-record-kernel-signature}\\
\mathbb P_{\theta,V}&:(\mathsf X,\mathscr X)\longrightarrow\mathcal P\bigl(\mathsf Y\times\mathsf O,\mathscr Y\otimes\mathscr O\bigr),
&X&\longmapsto\mathbb P_{\theta,V}(\cdot\given X),
\label{eq:prob-generative-signature}\\
\mathbb Q_V&:(\mathsf O\times\mathsf X,\mathscr O\otimes\mathscr X)\longrightarrow\mathcal P\bigl(\mathsf Y,\mathscr Y\bigr),
&(o,X)&\longmapsto\mathbb Q_{V,o,X},
\label{eq:prob-recognition-signature}
\end{align}
where each record coordinate occurs once, $\mathbb P_{\theta,V}$ is constructed in \Cref{ch:generative}, and
$\mathbb Q_{V,o,X}$ is the coupling selected in \Cref{sec:prob-marginals}; neither is primitive data. The
identity access map is the unrestricted interface; a stricter one restricts what the local law may depend
on, not the population law (\Cref{prop:prob-access-restriction}). An agent's generative-model structure is
the pair $(q_i^m,\operatorname{ev}_i)$ of its model law and the evaluator of \Cref{def:prob-model-evaluation}.
The generative kernel reads only $\theta$, $X$, and latent coordinates (\Cref{req:gen-typing-prohibition}),
and a configuration prior, when declared, is a law on $(\mathsf X,\mathscr X)$.

\definitionheading{Model presentations and evaluated kernels}{def:prob-model-evaluation}
For each $i\in V$ and structural value $X$, let $(\mathsf S_i^X,\mathscr S_i^X)$ and
$(\mathsf T_i^X,\mathscr T_i^X)$ contain every source and output coordinate of agent $i$'s generative
factor, and let $(\mathsf M_i,\mathscr M_i)$ be its presentation space, so that a model sample is
$m_i\in\mathsf M_i$ and the model law is $q_i^m\in\mathcal P(\mathsf M_i)$. The evaluator is a map
\begin{equation}
\operatorname{ev}_i:\mathsf M_i\longrightarrow\operatorname{Kern}(\mathsf S_i^X,\mathsf T_i^X),
\qquad
m_i\longmapsto K^X_{i,m_i},
\label{eq:prob-model-evaluation}
\end{equation}
with $(m_i,z)\mapsto K^X_{i,m_i}(z,B)$ measurable for every $B\in\mathscr T_i^X$ and
$K^X_{i,m_i}(z,\mathsf T_i^X)=1$. Each source space retains every observation, structural, or parent
coordinate its factor requires. In the definite case $q_i^m=\delta_{m_i}$. Neither a model law, a
recognition parameter, nor a posterior is an argument of an evaluated kernel.

\section{Reference measures}
\label{sec:prob-reference}

\definitionheading{Reference measures}{def:prob-reference-measures}
Each Euclidean coordinate carries its Borel sigma-algebra and Lebesgue measure, and each countable
coordinate its power set and counting measure. A mixed real coordinate carries one parameter-independent
reference for the whole family of laws placed on it, generative and recognition alike: when the union
$A_*$ of the atom sets of that family is countable, it carries $\nu_{\mathrm{mix}}$ of
\eqref{eq:prob-mixed-family-reference}; when $A_*$ is uncountable no density model is declared on that
coordinate (\Cref{thm:prob-mixed-family-common-domination}). Any other fiber carries an explicitly
declared sigma-finite measure. With $\nu_i^k,\nu_i^m,\nu_a$ the component measures, the product
references are
\begin{equation}
\nu^{\mathsf Y}=\bigotimes_{i\in V}\bigl(\nu_i^k\otimes\nu_i^m\bigr),
\qquad
\nu^{\mathsf O}=\bigotimes_{a\in\mathfrak A}\nu_a.
\label{eq:prob-reference-measures}
\end{equation}

\theoremheading{Common domination for a mixed family}{thm:prob-mixed-family-common-domination}
Let $(\mu_t)_{t\in\mathbb T}$ be any family of Borel probability measures on $\R$ with no
singular-continuous part, written $\mu_t=f_t\operatorname{Leb}+\sum_{a\in A_t}p_t(a)\delta_a$ with
$A_t:=\{a\in\R:\mu_t(\{a\})>0\}$. A single sigma-finite Borel measure dominates every $\mu_t$ if and only
if $A_*:=\bigcup_{t\in\mathbb T}A_t$ is countable. When $A_*=\{a_1,a_2,\ldots\}$, one such
parameter-independent measure is
\begin{equation}
\nu_{\mathrm{mix}}:=\operatorname{Leb}+\sum_{n\geq1}2^{-n}\delta_{a_n}.
\label{eq:prob-mixed-family-reference}
\end{equation}
In particular, the moving-atom family $\mu_t=\tfrac12\mathcal N(0,1)+\tfrac12\delta_t$, $t\in\R$, has no
common sigma-finite dominating measure. \status{ESTABLISHED}

The proof is in \Cref{app:measure-lemmas}.

\remarkheading{Dirac families are not dominated}{sec:exp-dominated-nondominated}
On an uncountable standard Borel space, an uncountable family of Dirac laws admits no common
sigma-finite dominating measure, so density charts cannot be inferred from a general law fiber.
\status{ESTABLISHED}

\noindent\emph{Proof.}
A dominating measure must give positive mass to every atom of the family, and a sigma-finite measure
has at most countably many points of positive mass: on each finite-measure member of a countable cover,
the points of mass at least $1/n$ are finitely many. \hfill$\square$

\section{Kernels first, densities second}
\label{sec:prob-kernels}

\definitionheading{Normalized kernels; densities as derivatives}{def:prob-normalized-kernels}
The objects \eqref{eq:prob-generative-signature} and \eqref{eq:prob-recognition-signature} are
probability kernels. When $\mathbb P_{\theta,V}(\cdot\given X)\ll\nu^{\mathsf Y}\otimes\nu^{\mathsf O}$,
its density, the Radon--Nikodym derivative $p_{\theta,V}(o,y\given X)$, and the evidence, its marginal
density, are determined by
\begin{align}
\mathbb P_{\theta,V}(dy,do\given X)&=p_{\theta,V}(o,y\given X)\nu^{\mathsf Y}(dy)\nu^{\mathsf O}(do),
\label{eq:prob-generative-density}\\
p_{\theta,V}(o\given X)&=\int_{\mathsf Y}p_{\theta,V}(o,y\given X)\nu^{\mathsf Y}(dy).
\label{eq:prob-evidence}
\end{align}
When $\mathbb Q_{V,o,X}\ll\nu^{\mathsf Y}$ its density is written $q_{V,o,X}(y)$. Densities are determined
up to null sets of the reference, and versions that agree outside such a set denote the same density.
Measurability in $(o,X)$ of the evidence and of recognition expectations is
\Cref{prop:prob-kernel-integration-measurability}, and jointly measurable densities are
\Cref{cor:prob-common-reference-joint-density}; an unnormalized energy in a later chapter comes with a finite
positive normalizer.

\hypothesisheading{Regular-conditional convention}{hyp:prob-regular-conditional-convention}
For fixed $(\theta,X)$ the observation marginal and the posterior, a regular conditional probability, are
\begin{align}
\mathbb P_{\theta,V,X}^{\mathsf O}(A)&:=\mathbb P_{\theta,V}(\mathsf Y\times A\given X),
&A&\in\mathscr O,
\label{eq:prob-observation-marginal}\\
\boldsymbol\Pi_{\theta,V,o,X}(B)&=\mathbb P_{\theta,V}(B\given o,X),
&B&\in\mathscr Y,
\label{eq:prob-rcp}
\end{align}
the posterior being unique for $\mathbb P_{\theta,V,X}^{\mathsf O}$-almost every $o$. When the domination of
\eqref{eq:prob-generative-density} holds, \eqref{eq:prob-evidence} is a density of
$\mathbb P_{\theta,V,X}^{\mathsf O}$, and after a change on a $\nu^{\mathsf O}$-null set there is a
measurable set $\mathsf O_{\theta,V,X}^{\mathrm{reg}}$ with
\begin{align}
&\mathbb P_{\theta,V,X}^{\mathsf O}\bigl(\mathsf O_{\theta,V,X}^{\mathrm{reg}}\bigr)=1,
\qquad
0<p_{\theta,V}(o\given X)<\infty
\quad\text{on }\mathsf O_{\theta,V,X}^{\mathrm{reg}},
\label{eq:prob-regular-observations}\\
&\boldsymbol\Pi_{\theta,V,o,X}(B)
=\frac{\int_Bp_{\theta,V}(o,y\given X)\nu^{\mathsf Y}(dy)}{p_{\theta,V}(o\given X)},
\qquad o\in\mathsf O_{\theta,V,X}^{\mathrm{reg}}.
\label{eq:prob-rcp-density}
\end{align}
Every fixed-observation posterior, evidence-bound, or free-energy statement is asserted for
$\mathbb P_{\theta,V,X}^{\mathsf O}$-almost every $o$, represented by
$o\in\mathsf O_{\theta,V,X}^{\mathrm{reg}}$. The hypothesis excludes everywhere-pointwise statements
unless a jointly measurable density, an evidence representative, and a conditional version are declared
as model data; the joint measure alone does not determine a posterior at an exceptional observation
(\Cref{rem:prob-version-witness}).

\openproblemheading{Nondominated likelihood families}{open:prob-nondominated-likelihood}
The density formulas here use, for every record kernel $K_a$, one sigma-finite $\nu_a$ dominating
$\{K_a(X,y_{\partial a},\cdot)\}$ and a jointly measurable density. For a family without one, such as
$K_a(X,y_{\partial a},\cdot)=\delta_{h(y_{\partial a})}$ with $h$ taking uncountably many values
(\Cref{sec:exp-dominated-nondominated}), the joint law and its regular conditional posterior exist, but the
likelihood $dK_a(X,y_{\partial a},\cdot)/d\nu_a(o_a)$ is undefined and the evidence density can fail to
exist. Open is a treatment of the evidence bound, its local disintegrations, and the coarse channels at the
level of joint measures and regular conditional laws, or hypotheses restoring a jointly measurable
likelihood, with a rule for propagating record-null choices of versions. It is settled by such a treatment
with proofs that the exact identities of \Cref{ch:elbo,ch:coarsegraining} hold in it, or by a record family
for which they provably have no measure-level counterpart.

\section{Absolute continuity}
\label{sec:prob-abscont}

\definitionheading{Logarithmic density ratios and relative entropy}{def:prob-log-density-ratio-convention}
Wherever a logarithmic density ratio or its expectation is written, the numerator measure is absolutely
continuous with respect to the denominator measure. Relative entropy is
\begin{equation}
\KL(Q\Vert P)=
\begin{cases}
\displaystyle\int\log\Bigl(\frac{dQ}{dP}(y)\Bigr)Q(dy), & Q\ll P,\\[1.2ex]
+\infty, & \text{otherwise},
\end{cases}
\label{eq:prob-kl-definition}
\end{equation}
the first branch being the divergence of \citet{Kullback1951} and the value $+\infty$ off absolute
continuity the convention of \citet{Csiszar1967}. The negative part of the integrand is $Q$-integrable,
since $\int_{\{dQ/dP<1\}}-\log(dQ/dP)dQ\leq1/e$, so the first branch lies in $[0,+\infty]$.

\propositionheading{Density criterion for absolute continuity}{prop:prob-density-absolute-continuity}
Let $Q$ and $P$ be probability measures with $Q\ll\nu$ and $P\ll\nu$ for a sigma-finite $\nu$, with
densities $q$ and $p$, and put $Z=\{p=0\}$. Then
\begin{align}
&Q\ll P
\iff
Q(Z)=0
\iff
q=0\quad\nu\text{-almost everywhere on }Z,
\label{eq:prob-abscont-criterion}\\
&Q(S)=1,
\qquad
P(S)=0,
\label{eq:prob-direct-null-set}
\end{align}
where the first line holds, and whether or not $Q$ is dominated by $\nu$, a measurable $S$ satisfying the
second gives $Q\not\ll P$ and $\KL(Q\Vert P)=+\infty$. \status{ESTABLISHED}

\noindent\emph{Proof.}
First $P(Z)=\int_Zp(y)\nu(dy)=0$, so $Q\ll P$ gives $Q(Z)=0$. Conversely, if $Q(Z)=0$ and $P(A)=0$, then
$\nu(A\cap\{p>0\})=0$, so $Q(A\cap\{p>0\})=0$, while $Q(A\cap Z)\leq Q(Z)=0$. The second equivalence
follows from $Q(Z)=\int_Zq(y)\nu(dy)$. The last claim is the definition of absolute continuity together with
\eqref{eq:prob-kl-definition}. \hfill$\square$

A law carried by a Lebesgue-null set, such as a proper affine subspace, has no Lebesgue density, and
\eqref{eq:prob-direct-null-set} shows its singularity (\Cref{prop:elbo-subspace-support-singular}).

\section{The relative-entropy chain rule and data processing}
\label{sec:prob-kl-chain-rule}

\lemmaheading{Relative-entropy chain rule, data processing, and bijection invariance}{lem:prob-kl-chain-rule}
Let $P,Q$ be laws on $\mathsf Y_1\times\mathsf Y_2$ with first marginals $P_1,Q_1$ and conditional kernels
$P_{2|1},Q_{2|1}$. Then:
\begin{enumerate}
\item[(a)] the map $y_1\mapsto\KL(Q_{2|1}(y_1,\cdot)\Vert P_{2|1}(y_1,\cdot))$ is measurable and
\begin{equation}
\KL(Q\Vert P)
=\KL(Q_1\Vert P_1)
+\int\KL\bigl(Q_{2|1}(y_1,\cdot)\big\Vert P_{2|1}(y_1,\cdot)\bigr)Q_1(dy_1)
\label{eq:elbo-kl-chain-rule}
\end{equation}
in $[0,+\infty]$, with a right side that does not depend on the chosen conditional versions;
\item[(b)] on a finite product $\mathsf Y=\prod_{b=1}^B\mathsf Y_b$, for any law $Q$ with marginals $Q_b$
and any laws $\rho_b$ on $\mathsf Y_b$,
\begin{equation}
\KL\Bigl(Q\Big\Vert\bigotimes_b\rho_b\Bigr)
=\KL\Bigl(Q\Big\Vert\bigotimes_bQ_b\Bigr)+\sum_b\KL(Q_b\Vert\rho_b)
\label{eq:prob-kl-product-reference}
\end{equation}
in $[0,+\infty]$;
\item[(c)] for a law $\mu$ on $\mathsf Y_1$, the minimum of $\KL(Q\Vert P)$ over all laws $Q$ with
$Q_1=\mu$ is $\KL(\mu\Vert P_1)$, attained by $\mu(dy_1)P_{2|1}(y_1,dy_2)$, uniquely when
$\KL(\mu\Vert P_1)<\infty$;
\item[(d)]\label{thm:cg-kl-dpi-extended} for laws $P,Q$ on $\mathsf Y$ and every Markov kernel
$K:\mathsf Y\rightsquigarrow\mathsf Z$,
\begin{equation}
\KL(QK\Vert PK)\leq\KL(Q\Vert P)
\qquad\text{in }[0,+\infty];
\label{eq:cg-kl-dpi-extended}
\end{equation}
\item[(e)] for a bimeasurable bijection $A:\mathsf Y\to\mathsf Y'$,
\begin{align}
\frac{d(A_\#Q)}{d(A_\#P)}&=\frac{dQ}{dP}\circ A^{-1}
\quad A_\#P\text{-almost surely, if }Q\ll P,
\label{eq:cg-rn-bimeasurable}\\
\KL(A_\#Q\Vert A_\#P)&=\KL(Q\Vert P)
\quad\text{in }[0,+\infty].
\label{eq:cg-kl-bimeasurable-invariance}
\end{align}
\end{enumerate}
\status{ESTABLISHED}

The proof is in \Cref{app:measure-lemmas}. Clauses (a), (b), and (d) need no absolute-continuity or
finiteness hypothesis.

\section{Recognition marginals, couplings, and selection}
\label{sec:prob-marginals}

\definitionheading{Local recognition laws and their coupling class}{def:prob-recognition-marginals}
With $\operatorname{pr}_i:\mathsf Y\to\mathsf Y_i$ the agent-block projection and
$\operatorname{pr}_i^k,\operatorname{pr}_i^m$ the coordinate projections of $\mathsf Y_i$, for each admitted
$(o,X)$ put
\begin{align}
\operatorname{Cpl}(\mathbb Q_{\bullet,o,X})
&:=\Bigl\{\mathbb J\in\mathcal P(\mathsf Y):(\operatorname{pr}_i)_\#\mathbb J=\mathbb Q_{i,o,X}\text{ for every }i\in V\Bigr\},
\label{eq:prob-recognition-coupling-class}\\
q_i^{b;o,X}&:=(\operatorname{pr}_i^k)_\#\mathbb Q_{i,o,X},
\qquad
q_i^{m;o,X}:=(\operatorname{pr}_i^m)_\#\mathbb Q_{i,o,X}.
\label{eq:prob-recognition-marginals}
\end{align}
The belief and model recognition marginals $q_i^{b;o,X},q_i^{m;o,X}$ are the values at $c_*$, in the frame
$u_i$, of agent $i$'s sections $q_i^b,q_i^m$ at the admitted $(o,X)$. They are local data; a member of
\eqref{eq:prob-recognition-coupling-class} is population-level dependence data.

\propositionheading{Existence and convexity of the coupling class}{prop:prob-coupling-existence-convexity}
For every admitted $(o,X)$ the class $\operatorname{Cpl}(\mathbb Q_{\bullet,o,X})$ is nonempty and convex,
and
\begin{equation}
(o,X)\longmapsto\bigotimes_{i\in V}\mathbb Q_{i,o,X}
\label{eq:prob-product-coupling-kernel}
\end{equation}
is a measurable product-coupling kernel. \status{ESTABLISHED}

\noindent\emph{Proof.}
On measurable rectangles the product is a product of measurable functions, and a monotone-class argument
extends measurability; its $i$-th marginal is $\mathbb Q_{i,o,X}$ because the other factors have unit
mass. Marginalization is affine, so convex combinations of couplings are couplings. \hfill$\square$

\definitionheading{Selected population recognition kernel}{hyp:prob-measurable-coupling-selector}
\label{def:elbo-recognition-kernel}
A selection rule $\mathsf S$ is declared, with law spaces carrying their evaluation sigma-algebras, such
that the population recognition kernel
\begin{align}
\mathbb Q_V&:(\mathsf O\times\mathsf X,\mathscr O\otimes\mathscr X)\longrightarrow\mathcal P(\mathsf Y,\mathscr Y),
&(o,X)&\longmapsto\mathbb Q_{V,o,X},
\label{eq:elbo-recognition-kernel}\\
\mathbb Q_{V,o,X}&=\mathsf S_{o,X}\bigl((\mathbb Q_{i,o,X})_{i\in V}\bigr)\in\operatorname{Cpl}(\mathbb Q_{\bullet,o,X}),
\label{eq:prob-selected-recognition-coupling}
\end{align}
has $(o,X)\mapsto\mathbb Q_{V,o,X}(B)$ measurable for every $B\in\mathscr Y$. The law may correlate agents
and channels; the product coupling is one admissible selection, not the default, and copulas, common latent
variables, ordered conditionals, or graphical families supply others.

\hypothesisheading{Constrained coupling selection}{hyp:prob-constrained-coupling-selection}
When the selection must lie in an imposed subfamily
$\operatorname{Cpl}^{\mathrm{res}}(\mathbb Q_{\bullet,o,X})\subseteq\operatorname{Cpl}(\mathbb Q_{\bullet,o,X})$,
every fiber $\operatorname{Cpl}^{\mathrm{res}}(\mathbb Q_{\bullet,o,X})$ is assumed nonempty,
\label{hyp:prob-constrained-coupling-feasibility}
and a selector $\mathsf S$ from these fibers is assumed with
$(o,X)\mapsto\mathsf S_{o,X}((\mathbb Q_{i,o,X})_{i\in V})(B)$ measurable for every $B\in\mathscr Y$. The
hypothesis excludes restricted families whose fibers are empty or admit no measurable selection; the two
conditions are independent, and the product coupling proves neither.

\propositionheading{Coordinate marginals do not determine a sufficiently rich joint}{prop:prob-marginals-do-not-determine-joint}
For a finite family of standard Borel blocks, the class $\operatorname{Cpl}((\mu_i)_{i\in V})$ is a
singleton if and only if at most one marginal $\mu_i$ is non-Dirac. Hence, when two local recognition
laws are non-Dirac, the local marginals do not determine $\mathbb Q_{V,o,X}$. \status{ESTABLISHED}

\noindent\emph{Proof.}
Let $\mu_1,\mu_2$ be non-Dirac. A non-Dirac law charges some measurable set with probability in $(0,1)$,
since otherwise the members or complements of probability one of a countable separating generator would
meet in a set of probability one with at most one point. Choose such sets $A,B$ and, for
$0<|\varepsilon|<1$, put
\begin{equation}
d\mathbb J_\varepsilon(x,x')
=\bigl[1+\varepsilon(\mathbf 1_A(x)-\mu_1(A))(\mathbf 1_B(x')-\mu_2(B))\bigr]d\mu_1(x)d\mu_2(x').
\label{eq:prob-coupling-perturbation}
\end{equation}
The bracket is at least $1-|\varepsilon|>0$ and integrating in either variable removes the perturbation,
so $\mathbb J_\varepsilon$ couples $\mu_1,\mu_2$, while
$\mathbb J_\varepsilon(A\times B)=\mu_1(A)\mu_2(B)+\varepsilon\mu_1(A)(1-\mu_1(A))\mu_2(B)(1-\mu_2(B))$
varies with $\varepsilon$; tensoring with the other marginals gives distinct couplings of the family.
Conversely, if at most one marginal is non-Dirac, every coupling is carried by the Dirac atoms in the
other coordinates, and its law there is the remaining marginal. \hfill$\square$

\propositionheading{Limits of access restriction}{prop:prob-access-restriction}
A restriction of the access maps constrains the local laws and not the population law, and it has an
evidence price that no selector removes.
\begin{enumerate}
\item[(i)] Let $V=\{1,2\}$ with binary blocks, let both local laws be $\operatorname{Ber}(\tfrac12)$ at every
observation, so that every access map is respected, and fix $0<r<1$. The selector that returns the
coupling with $\mathbb Q_{V,o,X}(y_1=y_2)=(1+r)/2$ when a record $o_2$ of agent $2$ equals $1$ and
$(1-r)/2$ otherwise satisfies \eqref{eq:prob-selected-recognition-coupling}, and the population law
depends on $o_2$ while no local law does. A decentralized information claim therefore also requires a
selector indexed by $((\operatorname{Acc}_i(o,X))_{i\in V},X)$ alone.
\item[(ii)] Let $V=\{1,2\}$ with binary latents, the prior with mass $0.45$ on each of $(0,0)$ and $(1,1)$
and $0.05$ on each of $(0,1)$ and $(1,0)$, and singleton records $o_1,o_2$ that report $y_1,y_2$
correctly with probability $0.8$, independently given $y$. Let $\operatorname{Acc}_1(o,X)=o_1$, so that
$\mathbb Q_{1,o,X}=\operatorname{Ber}(t)$ is the same at $o=(1,0)$ and $o=(1,1)$. The posterior marginals
of $y_1$ at these observations, written as the probabilities of $y_1=0$ and $y_1=1$, are
$\Pi_1^{(0)}=(37/89,52/89)$ and $\Pi_1^{(1)}=(13/161,148/161)$, and
every recognition law with this access restriction satisfies
\begin{equation}
\begin{aligned}
\max_{o_2\in\{0,1\}}\KL\bigl(\mathbb Q_{V,(1,o_2),X}\big\Vert\boldsymbol\Pi_{\theta,V,(1,o_2),X}\bigr)
&\geq\max_{o_2\in\{0,1\}}\KL\bigl(\operatorname{Ber}(t)\big\Vert\Pi_1^{(o_2)}\bigr)\\
&\geq\sup_{\xi\in[0,1]}\bigl(-\log Z_\xi\bigr)=0.08849\ldots,
\end{aligned}
\label{eq:prob-access-gap-floor}
\end{equation}
where $Z_\xi=\sum_{y_1}\Pi_1^{(0)}(y_1)^\xi\Pi_1^{(1)}(y_1)^{1-\xi}$. The last bound is sharp: the supremum
is attained at $\xi\approx0.548$, and at $t\approx0.783$ the middle term equals it. Hence the equality
condition of \Cref{cor:elbo-bound-tightness} lies outside the access-restricted class, and the exact
E-coordinate of \Cref{prop:elbo-exact-e-coordinate}, which optimizes over every law with a prescribed
complement marginal, need not be access restricted.
\end{enumerate}
\status{ESTABLISHED}

\noindent\emph{Proof.}
(i) Both couplings have $\operatorname{Ber}(\tfrac12)$ marginals. (ii) Bayes' rule gives the posterior
marginals, and projection onto $y_1$ with \Cref{lem:prob-kl-chain-rule}(d) gives the first inequality. For
$\xi\in[0,1]$ let $\Pi^{[\xi]}=Z_\xi^{-1}(\Pi_1^{(0)})^\xi(\Pi_1^{(1)})^{1-\xi}$. Expanding logarithms gives
$\xi\KL(\operatorname{Ber}(t)\Vert\Pi_1^{(0)})+(1-\xi)\KL(\operatorname{Ber}(t)\Vert\Pi_1^{(1)})
=\KL(\operatorname{Ber}(t)\Vert\Pi^{[\xi]})-\log Z_\xi\geq-\log Z_\xi$, and a maximum dominates every convex
combination. The concave function $\xi\mapsto-\log Z_\xi$ has derivative
$\E_{\Pi^{[\xi]}}\log(\Pi_1^{(1)}/\Pi_1^{(0)})$, which vanishes at its interior maximizer $\xi_\star$, where
the two divergences from $\Pi^{[\xi_\star]}$ are therefore equal; the identity with
$\operatorname{Ber}(t)=\Pi^{[\xi_\star]}$ shows that both equal $-\log Z_{\xi_\star}$, the Chernoff
information \citep[Section~11.9]{CoverThomas2006}. The value and maximizers follow from the closed forms,
and every access-restricted law has a positive gap at one of the two observations. \hfill$\square$

\section{Levels of description}
\label{sec:prob-levels}

\definitionheading{Sample, law-point, and configuration levels}{def:prob-description-levels}
At the sample level the agent block is $\mathsf Y_i$ of \Cref{def:prob-product-sample-spaces}, the
population recognition law $\mathbb Q_{V,o,X}$ is a law of samples, and the belief of agent $i$ is the
recognition marginal $q_i^{b;o,X}$ of \eqref{eq:prob-recognition-marginals}. Generative kernels and the
record kernels \eqref{eq:prob-record-kernel-signature} read sample coordinates. At the law-point level, law
points are latent coordinates, and the law-point blocks are
\begin{equation}
\mathsf Y^{\mathrm{law}}_i=(\mathcal E_b)_{c_*}\times(\mathcal E_m)_{c_*}\ni\bigl(\mathfrak q_i^b,\mathfrak q_i^m\bigr),
\qquad
\mathsf Y^{\mathrm{law}}=\prod_{i\in V}\mathsf Y^{\mathrm{law}}_i,
\label{eq:prob-law-point-space}
\end{equation}
possibly augmented by further latent coordinates; through the frame $u_i$, $\mathfrak q_i^b$ and
$\mathfrak q_i^m$ are laws on $\mathsf K_i$ and $\mathsf M_i$. A generative law at this level is a law on
$\mathsf Y^{\mathrm{law}}\times\mathsf O$ whose kernels may read law points, and a recognition law is a law
over law points. At the configuration level the latent state is
$(w,y)\in\mathsf Y^{\mathrm{cfg}}\times\mathsf Y$, where the configuration space $\mathsf Y^{\mathrm{cfg}}$
carries a measurable presentation map
$\operatorname{pres}:\mathsf Y^{\mathrm{cfg}}\to\mathsf Y^{\mathrm{law}}$ and may carry further latent
coordinates such as random links. A lifted generative law is a law $\mathbb P^{\mathrm{cfg}}_{\theta,V}(dw,dy,do\given X)$, and a lifted
recognition law has the form $\mathbb J(dw,dy)=\mathbb J^{\mathrm{cfg}}(dw)\kappa^{\mathrm{pres}}(w,dy)$ with a
presentation kernel $\kappa^{\mathrm{pres}}:\mathsf Y^{\mathrm{cfg}}\rightsquigarrow\mathsf Y$ whose marginal
at each agent is the law presented by the corresponding component of $\operatorname{pres}(w)$.

The levels define different optimization problems, so an energy enters a free energy only after its level
is named. The source energies of \eqref{eq:obs-attention-likelihood} and every record kernel are
sample-level. The transported divergences \eqref{eq:full-graph-transported-divergences} and the role
divergences \eqref{eq:obs-selective-role-divergence} are measurable $[0,+\infty]$-valued functions of law
points (\Cref{cor:prob-kernel-kl-measurability}), admissible as generative energies when their arguments
are latent law points; read at the sample level they would make a generative factor read recognition
marginals, which \Cref{req:gen-typing-prohibition} excludes. A block admitted by \Cref{hyp:local-interaction-kernels} with
law points is a law-point or configuration block, and the formation model
\eqref{eq:rg-formation-configuration} is configuration-level, with records that read the configuration.

\propositionheading{A law-point energy is not a sample energy}{prop:prob-law-point-energy}
Let $q_i$ be a receiver law and $p_1,p_2$ source laws in $\mathcal B_x$ with
$\KL(q_i\Vert\widehat\rho_x(\Theta_e^x)p_k)<\infty$ for $k=1,2$, whose transported laws have densities
against a common dominating measure that differ on a set of positive $q_i$-measure. Then for $0<s<1$
\[
\KL\bigl(q_i\Vert\widehat\rho_x(\Theta_e^x)(sp_1+(1-s)p_2)\bigr)
<s\KL\bigl(q_i\Vert\widehat\rho_x(\Theta_e^x)p_1\bigr)+(1-s)\KL\bigl(q_i\Vert\widehat\rho_x(\Theta_e^x)p_2\bigr),
\]
and no recognition-independent sample energy $D$ satisfies
$\E_{q_i\otimes q}D=\KL(q_i\Vert\widehat\rho_x(\Theta_e^x)q)$ at $q=p_1$, at $q=p_2$, and at one strict
mixture of them. \status{ESTABLISHED}

\noindent\emph{Proof.}
Pushforward preserves mixtures. Let $\alpha,\beta$ be the $q_i$-densities of the absolutely continuous parts
of the two transported laws. Finite divergences give $q_i\ll$ each transported law, so $\alpha,\beta>0$
$q_i$-almost everywhere and the divergences are $\int-\log\alpha(y)q_i(dy)$ and $\int-\log\beta(y)q_i(dy)$,
well defined because $\int(\log\alpha(y))^+q_i(dy)\leq\int\alpha(y)q_i(dy)\leq1$. The mixture has relative
density $s\alpha+(1-s)\beta$, and $-\log(s\alpha+(1-s)\beta)\leq-s\log\alpha-(1-s)\log\beta$, strictly on the
set $\{\alpha\neq\beta\}$ of positive $q_i$-measure; integrating gives the strict inequality. An energy $D$
contributes $\E_{q_i\otimes q}D$, affine in $q$ along mixtures, so matching the divergence at $p_1$ and $p_2$
forces the affine value at the mixture, which the strict inequality excludes. \hfill$\square$

\propositionheading{Point-mass recognition of law points has infinite cost}{prop:prob-law-point-point-mass}
At the law-point level let $\mathbb P^{\mathrm{law}}_{\theta,V}(\cdot\given X)$ be a generative law with
latent marginal $\mathbb P^{\mathrm{law},\mathsf Y}_{\theta,V}$ and posterior
$\boldsymbol\Pi^{\mathrm{law}}_{\theta,V,o,X}\ll\mathbb P^{\mathrm{law},\mathsf Y}_{\theta,V}$ at an observation
with finite positive evidence; the domination holds when every record kernel has a density against a fixed
reference. If $\bar{\mathfrak q}$ is not an atom of the prior, then
\begin{equation}
\KL\bigl(\delta_{\bar{\mathfrak q}}\big\Vert\boldsymbol\Pi^{\mathrm{law}}_{\theta,V,o,X}\bigr)=+\infty,
\label{eq:prob-law-point-point-mass}
\end{equation}
and the extended evidence bound of $\delta_{\bar{\mathfrak q}}$ is $-\infty$. Consequently, against a
non-atomic law-point prior, a row or an update computed at fixed law points is not the recognition
optimum of a finite evidence bound at the law-point level: with the law points latent it is a generative
posterior quantity conditional on them, and with the law points declared as parameters in $(\theta,X)$ it
is a function of parameters whose change is an M-coordinate subject to
\Cref{def:elbo-generalized-m-acceptance}. \status{ESTABLISHED}

\noindent\emph{Proof.}
Singletons are measurable, and domination gives
$\boldsymbol\Pi^{\mathrm{law}}_{\theta,V,o,X}(\{\bar{\mathfrak q}\})=0$, so \eqref{eq:prob-direct-null-set} holds for
$\delta_{\bar{\mathfrak q}}$ with $S=\{\bar{\mathfrak q}\}$. The extended bound
\eqref{eq:elbo-extended-definition} subtracts this divergence from a finite log evidence, and a recognition
optimum with finite bound has finite divergence. \hfill$\square$

\propositionheading{Configuration-lift identity}{prop:prob-configuration-lift-identity}
Let $\kappa^{\mathrm{cfg}}:\mathsf Y\rightsquigarrow\mathsf Y^{\mathrm{cfg}}$ be a Markov kernel, and let the lifted
generative law be $\mathbb P_{\theta,V}(dy,do\given X)\kappa^{\mathrm{cfg}}(y,dw)$, so that its
$(y,o)$ marginal is the sample-level law and its records read only sample coordinates. Then at every
$o\in\mathsf O^{\mathrm{reg}}_{\theta,V,X}$ the evidence is unchanged, the lifted posterior may be selected
as $\boldsymbol\Pi_{\theta,V,o,X}(dy)\kappa^{\mathrm{cfg}}(y,dw)$, and every lifted recognition law $\mathbb J$
with sample marginal $\bar{\mathbb Q}$ and conditional kernel $\mathbb J(dw\given y)$ satisfies
\begin{equation}
\Lelbo^{\mathrm{ext,cfg}}_V(\mathbb J;X,o)
=\Lelbo^{\mathrm{ext}}_V(\bar{\mathbb Q};X,o)
-\E_{\bar{\mathbb Q}}\Bigl[\KL\bigl(\mathbb J(\cdot\given Y)\big\Vert\kappa^{\mathrm{cfg}}(Y,\cdot)\bigr)\Bigr]
\label{eq:prob-configuration-lift-identity}
\end{equation}
in $[-\infty,\log p_{\theta,V}(o\given X)]$, where $\Lelbo^{\mathrm{ext,cfg}}_V$ is
\eqref{eq:elbo-extended-definition} for the lifted law. The lift never raises the bound of the induced
sample law, and when that bound is finite it leaves it unchanged exactly when
$\mathbb J(\cdot\given y)=\kappa^{\mathrm{cfg}}(y,\cdot)$ for $\bar{\mathbb Q}$-almost every $y$.
\status{ESTABLISHED}

\noindent\emph{Proof.}
The $(y,o)$ marginal of the lifted law is $\mathbb P_{\theta,V}$, so the evidence is unchanged, and
disintegrating along $o$ gives the stated posterior. \Cref{lem:prob-kl-chain-rule}(a), with $y$ as the
first coordinate, gives the lifted divergence as $\KL(\bar{\mathbb Q}\Vert\boldsymbol\Pi_{\theta,V,o,X})$ plus
the conditional term in $[0,+\infty]$, and the conditional term vanishes exactly under the stated
almost-everywhere equality. Substitution in
\eqref{eq:elbo-extended-definition} subtracts two terms in $[0,+\infty]$ from a finite log evidence.
\hfill$\square$

\openproblemheading{A normalized law-point prior}{open:prob-law-point-prior}
No chapter constructs a normalized generative law on $\mathsf Y^{\mathrm{law}}$ for a general class of law
points, which the law-point energies \eqref{eq:full-graph-transported-divergences} and
\eqref{eq:obs-selective-role-divergence} need to be evidence-bound terms rather than a separately declared
potential. The law fibers supply no canonical sigma-finite reference ($\mathcal P(\mathsf K)$ is in general
infinite dimensional), and the Haar measure of a noncompact $G$ is infinite, so a Haar prior on a free orbit
is not normalizable; the formation model supplies one normalized configuration-level prior,
\eqref{eq:rg-formation-physical-prior}, for finite-dimensional Gaussian presentations with a proper,
non-invariant reference. It is settled for a class of law points by a sigma-finite reference on the law-point
or presentation space, a finite normalizer for the exponentiated energies, the transformation law of the prior
under the active gauge action, record kernels that read law points, and a compatibility relation between the
law-point posterior and the sample-level recognition marginals.

\section*{Physicist's summary}

Everything probabilistic at $c_*$ is a law or a kernel, with densities against references fixed per family,
and one chain rule in $[0,+\infty]$ does all later bookkeeping. Local recognition laws do not determine their
coupling, restricted access has an exact evidence price, and attention's transported divergence is a
law-point energy.
